\documentclass[10pt,a4paper]{amsart}
\usepackage[margin=19mm]{geometry}
\usepackage{amsmath,amssymb,mathtools}
\usepackage[scr=rsfs,cal=euler]{mathalfa}
\usepackage{xfrac}
\usepackage[unicode,hidelinks,bookmarks=false]{hyperref}
\newcommand{\ie}{i.e.\ }
\newcommand{\cf}{cf.\ }
\newcommand{\resp}{respectively\ }
\newcommand{\Subsection}{\subsection}
\providecommand{\nobreakdash}{\nobreak\hbox{-}}
\setlength{\emergencystretch}{2em}
\multlinegap0pt
\usepackage{amssymb}
\usepackage{xcolor}
\usepackage{braket}
\usepackage{mathtools}
\usepackage{tikz}
\newtheorem{theorem}{Theorem}
\newcommand{\bvec}[1]{\boldsymbol{#1}}
\newcommand{\symp}{\Omega}
\title{Symplectic Barnes-Wall GKP Codes:\\Deterministic $O(N \log^2 N)$ Decoding and Logarithmic Rate Scaling}
\author{Shanxiang Lyu}
\date{Source: arXiv:2608.00601v1 (CC BY 4.0)}
\begin{document}
\maketitle

\noindent\emph{Source and licence.} Symplectic Barnes-Wall GKP Codes:\\Deterministic $O(N \log^2 N)$ Decoding and Logarithmic Rate Scaling. Version arXiv:2608.00601v1 (CC BY 4.0), distributed by arXiv under the Creative Commons Attribution 4.0 International licence, http://creativecommons.org/licenses/by/4.0/, as recorded in the arXiv licence metadata for this exact version and shown on the abstract page https://arxiv.org/abs/2608.00601v1. Full licence notice: \texttt{licenses/arxiv-2608.00601-CC-BY-4.0.txt}.\par\smallskip

\noindent\textit{Editorial scope.} Counted unit: the article's theorem thm:main (source0.tex lines 170--176). The article's complete original proof of it is delivered with it (source0.tex lines 178--205, one \texttt{proof} environment of 28 lines). No mathematics is written, repaired or re-derived: the statement and the proof are the article's own lines, in the article's own notation, and no retained line is substituted or re-flowed.  Retained uncounted as the source-local context the proof needs: lines 161--168.  Nothing was omitted from any retained span, so the package carries no omission record.  No retained line contains a citation, so the wrapper carries no bibliography and no bibliography entry.  No retained line contains a figure environment, an image-inclusion command or drawing code, so no image is delivered and the package declares no asset dependency.  Editorial formatting only, in addition: an AMS \texttt{amsart} preamble that restates the article's own declarations and macros used by the retained lines, namely the environments \texttt{theorem} on separate counters, together with the standard packages those lines need.  The retained environments print in this document as Theorem 1 in order of appearance, whereas the article prints the same items with the numbering of its own full document.  The complete native source of the article as distributed in the arXiv e-print for this exact version is bundled unchanged as \texttt{sources/206554/source0.tex}.
We recursively define the symplectic Barnes-Wall generator. Let $m \geq 1$ index the recursion level, acting on $N = 2^{m-1}$ modes ($2^m$ phase-space dimensions):
%
\begin{align}
    G_1 &= I_2, \label{eq:g1} \\
    G_{m+1} &= \begin{pmatrix} G_m & 0 \\ G_m & R_m G_m \end{pmatrix}, \label{eq:grec}
\end{align}
%
where $R_m = I_{2^m} + \symp_{2^m}$ satisfies $R_m^T R_m = 2 I_{2^m}$. Note that $R_m/\sqrt{2}$ is simultaneously orthogonal and symplectic with unit determinant; $R_m$ itself is a $\sqrt{2}$-scaled orthogonal symplectic matrix with $\det(R_m) = 2^{2^{m-1}}$.

\begin{theorem}[Symplectic Integrality]\label{thm:main}
For any $m \geq 1$, the generator $G_m$ is invertible, and the normalized overlap matrix
\begin{equation}
K_m := G_m^T \symp_{2^m} G_m
\end{equation}
is integer-valued, antisymmetric, and non-singular. Consequently, $\tilde{M}_m = \sqrt{2\pi}\, G_m$ defines a valid multimode GKP code.
\end{theorem}

\begin{proof}
By induction. For $m=1$, $G_1 = I_2$, yielding $K_1 = \symp_2$, which is integer and antisymmetric.

Assume $K_m$ is integer and antisymmetric. Since $G_1 = I_2$ is integral and $R_m = I + \symp$ has entries in $\{0, \pm 1\}$, the recursion~\eqref{eq:grec} preserves integrality: $G_m \in M_{2^m}(\mathbb{Z})$ for all $m$. Define $S_m = G_m^T G_m$, which is symmetric and integral by the same closure argument.

Evaluating $K_{m+1}$ via block multiplication under the interleaved convention (so that $\symp_{2^{m+1}} = \symp_{2^m} \oplus \symp_{2^m}$):
\begin{align}
K_{m+1} &= \begin{pmatrix} G_m^T & G_m^T \\ 0 & G_m^T R_m^T \end{pmatrix}
           \begin{pmatrix} \symp_{2^m} & 0 \\ 0 & \symp_{2^m} \end{pmatrix}
           \begin{pmatrix} G_m & 0 \\ G_m & R_m G_m \end{pmatrix}.
\end{align}
Computing the four blocks:
\begin{align}
\text{Top-left:}&\quad G_m^T \symp G_m + G_m^T \symp G_m = 2K_m, \nonumber\\
\text{Top-right:}&\quad G_m^T \symp R_m G_m = G_m^T(\symp - I)G_m = K_m - S_m, \nonumber\\
\text{Bottom-left:}&\quad G_m^T R_m^T \symp G_m = G_m^T(\symp + I)G_m = K_m + S_m, \nonumber\\
\text{Bottom-right:}&\quad G_m^T R_m^T \symp R_m G_m = G_m^T (2\symp) G_m = 2K_m, \nonumber
\end{align}
where we used $\symp R_m = \symp(I+\symp) = \symp - I$, $R_m^T \symp = (I-\symp)\symp = \symp + I$, and $R_m^T \symp R_m = 2\symp$, all following from $\symp^2 = -I$, $\symp^T = -\symp$.

Thus:
\begin{equation}\label{eq:k-recursive}
K_{m+1} = \begin{pmatrix} 2K_m & K_m - S_m \\ K_m + S_m & 2K_m \end{pmatrix}.
\end{equation}
Antisymmetry holds because $(K_m - S_m)^T = K_m^T - S_m^T = -K_m - S_m = -(K_m + S_m)$. Integrality follows from closure under addition and multiplication over $\mathbb{Z}$. Non-singularity follows from $\det(G_{m+1}) = (\det G_m)^2 \det(R_m) \neq 0$.

By the Weyl commutation relation, $\hat{D}(\bvec{u})\hat{D}(\bvec{v}) = e^{-i \bvec{u}^T \symp \bvec{v}} \hat{D}(\bvec{v})\hat{D}(\bvec{u})$. Since every entry of $2\pi K_m = \tilde{M}_m^T \symp \tilde{M}_m$ is an integer multiple of $2\pi$, all generator displacements commute, yielding an abelian stabilizer group.
\end{proof}

\end{document}
